\section{Tinh bột và cellulose}

\subsection{Trạng thái tự nhiên và tính chất vật lí}

Tinh bột có nhiều trong hạt và củ như gạo, ngô, khoai, sắn. Đây là chất rắn dạng bột, thường có màu trắng, không tan trong nước lạnh. Khi đun nóng với nước, tinh bột trương nở và tạo thành hồ tinh bột.

Cellulose có nhiều trong thành tế bào thực vật, bông, gỗ và rơm rạ. Cellulose là chất rắn dạng sợi, không tan trong nước. Bông gần như chứa cellulose tinh khiết; gỗ và rơm rạ còn chứa nhiều thành phần khác.

Tinh bột và cellulose đều là carbohydrate có phân tử lớn, thường được biểu diễn bằng công thức dạng:
\[
\ce{(C6H10O5)_n}.
\]
Tuy có cùng công thức dạng mắt xích, tinh bột và cellulose có cấu tạo khác nhau nên không có cùng tính chất.

\begin{center}
\begin{tabularx}{0.96\textwidth}{>{\bfseries}p{2.5cm}YY}
\toprule
Đặc điểm & Tinh bột & Cellulose \\
\midrule
Nguồn tự nhiên & Gạo, ngô, khoai, sắn, hạt & Bông, gỗ, rơm rạ, thành tế bào thực vật \\
Trạng thái thường gặp & Chất rắn dạng bột & Chất rắn dạng sợi \\
Tính tan & Không tan trong nước lạnh; tạo hồ khi đun nóng với nước & Không tan trong nước \\
Vai trò chính & Dự trữ năng lượng ở thực vật & Tạo nên thành tế bào thực vật \\
Ứng dụng & Lương thực, thực phẩm; nguyên liệu sản xuất glucose và ethanol & Sản xuất giấy, vải sợi và một số nguyên liệu công nghiệp \\
\bottomrule
\end{tabularx}
\end{center}

\subsection{Tính chất hóa học của tinh bột}

\subsubsection*{Phản ứng màu với iodine}

Hồ tinh bột tác dụng với dung dịch iodine tạo màu xanh tím đặc trưng. Đây là phản ứng thường dùng để nhận biết tinh bột.

Khi đun nóng, màu xanh tím có thể mất đi; khi để nguội, màu xanh tím xuất hiện trở lại. Vì vậy, khi làm thí nghiệm cần quan sát mẫu sau khi thêm iodine và tránh nhầm sự mất màu tạm thời do đun nóng với việc tinh bột không có mặt.

\subsubsection*{Phản ứng thủy phân}

Khi đun nóng trong môi trường acid hoặc dưới tác dụng của enzyme thích hợp, tinh bột bị thủy phân. Sản phẩm cuối cùng là glucose:
\[
\ce{(C6H10O5)_n + nH2O ->[H^+,\ t^\circ] nC6H12O6}.
\]

Phản ứng này giải thích vì sao tinh bột có thể được dùng làm nguyên liệu sản xuất glucose. Trong cơ thể, enzyme cũng tham gia phân giải tinh bột trong quá trình tiêu hóa.

\subsection{Tính chất hóa học của cellulose}

Cellulose cũng có thể bị thủy phân trong điều kiện thích hợp, tạo glucose:
\[
\ce{(C6H10O5)_n + nH2O ->[H^+,\ t^\circ] nC6H12O6}.
\]

So với tinh bột, cellulose khó bị thủy phân hơn; quá trình chuyển hóa trong tự nhiên cần enzyme cellulase của một số vi sinh vật hoặc điều kiện xử lí thích hợp. Không nên suy ra rằng cellulose dễ bị thủy phân chỉ vì có cùng công thức dạng mắt xích với tinh bột.

Cellulose không tạo màu xanh tím đặc trưng với dung dịch iodine như hồ tinh bột.

\subsection{Sự tạo thành và vai trò trong cây xanh}

Trong quá trình quang hợp, cây xanh tạo ra glucose từ carbon dioxide và nước. Một phần glucose được chuyển hóa thành tinh bột để dự trữ; một phần được sử dụng để tạo cellulose, cấu thành thành tế bào thực vật.

Tinh bột là nguồn carbohydrate quan trọng trong lương thực. Cellulose là nguyên liệu thực vật phổ biến, được sử dụng trong sản xuất giấy, vải sợi và có thể được chuyển hóa thành glucose trong các quy trình xử lí phù hợp.

\begin{ghinho}{Nhận biết tinh bột}
\begin{itemize}
\item Hồ tinh bột + dung dịch iodine: xuất hiện màu xanh tím.
\item Glucose, saccharose và cellulose không cho màu xanh tím đặc trưng này trong phép thử trực tiếp.
\item Mẫu thực vật như chuối xanh, khoai, gạo có thể chứa tinh bột; khi thử iodine cần xét thành phần thực tế của mẫu.
\end{itemize}
\end{ghinho}

\begin{vidu}
Nhỏ vài giọt dung dịch iodine vào hai mẫu: hồ tinh bột và dung dịch glucose. Dự đoán hiện tượng và cho biết mẫu nào có thể nhận biết bằng thuốc thử này.
\end{vidu}
\begin{loigiai}
Hồ tinh bột chuyển sang màu xanh tím đặc trưng. Dung dịch glucose không xuất hiện màu xanh tím với iodine.

\textbf{Vậy,} có thể dùng dung dịch iodine để nhận biết hồ tinh bột trong hai mẫu.
\end{loigiai}
